\documentclass[11pt,a4paper]{article}

\usepackage[margin=2.2cm]{geometry}
\usepackage{booktabs}
\usepackage{array}
\usepackage{tabularx}
\usepackage{amsmath}
\usepackage{amssymb}
\usepackage[hidelinks]{hyperref}
\usepackage{enumitem}

\newcolumntype{L}{>{\raggedright\arraybackslash}p{0.16\textwidth}}
\newcolumntype{C}{>{\raggedright\arraybackslash}p{0.30\textwidth}}
\newcolumntype{R}{>{\raggedright\arraybackslash}p{0.44\textwidth}}

\newcommand{\code}[1]{\texttt{#1}}

\title{Pendulum Case:\\
A Comparison of \code{phys-vae} and the \code{interpretable-diffusion} PhysVAE Reimplementation\\
\large With a ranked diagnosis of pendulum-specific deviations, and what must change in PAID}
\author{Generated by Mistral Vibe}
\date{\today}

\begin{document}
\maketitle

\section*{Scope and sources}

This report is the pendulum counterpart of
\code{reports/advdif\_reimplementation\_comparison.tex} (advection--diffusion,
diagnosis D1--D8, since fixed on \code{main}). It compares the forced damped
pendulum experiment of the original \code{phys-vae} repository
(Takeishi \& Kalousis, \emph{Physics-Integrated Variational Autoencoders},
NeurIPS 2021; files \code{physvae/pendulum/model.py},
\code{physvae/pendulum/train.py}, \code{data/pendulum/generate.py},
\code{data/pendulum/makedata.sh}, \code{pendulum\_train.sh},
\code{pendulum\_cntfact.ipynb}) against the reimplementation in
\code{interpretable-diffusion} at \code{main} commit \code{a22f098}
(\code{src/models/phys\_vae.py}, \code{src/models/train\_phys\_vae.py},
\code{src/models/physics\_models.py}, \code{src/datasets/pivae.py},
\code{src/train.py}, \code{src/eval.py}, and
\code{configs/pendulum\_vae/26MAR13.yaml}), and then examines what the same
pendulum setup demands of PAID (\code{src/models/paid.py},
\code{src/models/train\_paid.py}, \code{configs/pendulum\_paid/*.yaml}).

Unlike the advdif case at the time of its report, the pendulum reimplementation
is in good shape: the D1--D8 fixes are inherited, the regularizer weights match
the original's \code{physnn} recipe exactly, and the model constructs and runs
(verified by a one-batch smoke test). The remaining deviations are ranked in
Section~\ref{sec:diagnosis}; the PAID-specific findings, led by a structural
problem with second-order ODE state initialization, are in
Section~\ref{sec:paid}. Two checks were executed in this environment (model
construction and one-batch loss/forward smoke tests, and a measurement of the
true angular velocity at PAID's condition frame); no training runs were
executed, and all accuracy statements are code-based.

\section{Setup at a glance}

\begin{table}[h]
\centering
\small
\begin{tabularx}{\textwidth}{LXX}
\toprule
 & \textbf{Original (\code{phys-vae})} & \textbf{Reimplementation} \\
\midrule
Framework & PyTorch + torchdiffeq & JAX / Flax nnx + diffrax \\
Task & $\ddot\theta = A\omega^2\cos(2\pi f t) - \gamma\dot\theta - \omega^2\sin\theta$; known physics = only $-\omega^2\sin\theta$; damping $\gamma$ and forcing $(A, f)$ absorbed by auxiliaries & Same concept: \code{pendulum\_model} implements $\dot\theta' = -\omega^2\sin\theta$ (\code{physics\_models.py:12--38}); damping and forcing absorbed by a \code{pre\_func} (in-ODE) and a \code{post\_func} (out-ODE) \\
Latents (\code{physnn} recipe) & $z_\mathrm{phy} = \omega$ (1-D), $z_{\mathrm{aux1}} \in \mathbb{R}^1$ (in-ODE), $z_{\mathrm{aux2}} \in \mathbb{R}^2$ (out-ODE) & Phys latent $=\omega$ (1-D), one pre-func with latent dim 1, one post-func with latent dim 2 (\code{26MAR13.yaml}) --- faithful \\
Data shape & $(N, 50)$ time series & $(N, 50, 1)$ with \code{retain\_dim: true}, flattened to $(B, 50)$ inside the model \\
\bottomrule
\end{tabularx}
\end{table}

\section{Data generation}
\label{sec:data}

\begin{table}[h]
\centering
\small
\begin{tabularx}{\textwidth}{LXX}
\toprule
\textbf{Aspect} & \textbf{Original} & \textbf{Reimplementation} \\
\midrule
ODE & $\ddot\theta = A\omega^2\cos(2\pi f t) - \gamma\dot\theta - \omega^2\sin\theta$ (\code{generate.py:11--15}) & Identical (\code{pivae.py:40--47}) \\
Parameter ranges & $\theta_0 \in [-1.57, 1.57]$, $\omega \in [0.785, 3.14]$, $\gamma \in [0, 0.8]$, $f \in [3.14, 6.28]$, $A \in [0, 40]$; $\dot\theta(0) = 0$ (\code{makedata.sh}) & Identical, including $\dot\theta(0) = 0$ (\code{pivae.py:113--152}) \\
Time grid & 50 frames, $dt = 0.05$, span $0 \ldots 2.45$; $dt$ stored in \code{args\_*.json} and read by the trainer & 50 frames, $dt = 2.5/50 = 0.05$, span $0 \ldots 2.45$ (\code{pivae.py:96--97}); the model derives $dt$ from the dataset config (\code{phys\_vae.py:595--609}) --- grids match, D1-class errors excluded \\
PDE solver & \code{scipy.solve\_ivp} DOP853, essentially exact & diffrax Euler at $dt_0 = 0.05$ (\code{pivae.py:143--153}) \\
Observation noise & $\sigma = 10^{-2}$ (\code{makedata.sh}) & \textbf{None} --- \code{generate\_pendulum\_data} returns the raw solver output (\code{pivae.py:155}) \\
Splits & train 1000 / valid 500 / test 1000, separate seeds & 3000 samples from one fixed seed, test 1000 / val 500 / train pool 1500 with 1000 drawn (\code{pivae.py:433--445}) \\
Saved targets & \code{init, omega, gamma, A, f} all saved (\code{generate.py:101--103}) & Only $\omega$ returned as the coefficient target (\code{pivae.py:155}) \\
\bottomrule
\end{tabularx}
\end{table}

\noindent Key observations:
\begin{itemize}[nosep]
  \item The Euler data generation is \emph{self-consistent} with the model's
        physics: the generator integrates the full ODE with Euler at
        $dt = 0.05$, and the model's physics term is also Euler at $dt = 0.05$
        (\code{integral\_level: 1}), so the decomposition
        $\dot y = f_\mathrm{phys} + f_\mathrm{aux}$ is exact in the Euler sense
        and the pre-function can in principle represent the missing damping and
        forcing exactly. The deviation from the original is data fidelity
        (Euler vs.\ DOP853 trajectories), not identifiability.
  \item The missing observation noise (N2) is the deviation with the largest
        interaction with the loss weighting, because the original's
        \code{x\_lnvar} $= -9$ encodes exactly that noise scale.
\end{itemize}

\section{Architectures}
\label{sec:arch}

\begin{table}[h]
\centering
\small
\begin{tabularx}{\textwidth}{LXX}
\toprule
\textbf{Component} & \textbf{Original} & \textbf{Reimplementation} \\
\midrule
Physics model & $\dot\theta' = -\omega^2\sin\theta$ with latent $\omega$ squared inside (\code{model.py:143--158}) & Identical (\code{physics\_models.py:34}) --- faithful \\
$\omega$ head & MLP + \textbf{Softplus} ($\omega > 0$ guaranteed), separate $\ln\sigma$ head; after sampling clamp to feasible $[0, 20]$ (\code{model.py:18, 101, 262--264}) & Linear $(\mu, \log\sigma)$ split; sampled latent hard-clipped to the model range $[0.392, 3.53]$ (\code{phys\_vae.py:748}) --- no positivity, dead gradients outside the range \\
Priors & $\omega$: mean $1.961$, $\ln\sigma^2 = 2\log(0.866 \cdot 3.138)$ (range $[0.392, 3.53]$, deliberately wider than the data range $[0.785, 3.14]$); auxiliaries $\mathcal{N}(0, 1)$ (\code{model.py:171--177}) & Identical formula and range (\code{phys\_vae.py:640--667}); auxiliaries $\mathcal{N}(0,1)$ broadcast --- faithful \\
Encoder & Per-latent feature extractor MLP $[50]{\to}[128,128]{\to}256$ (activation on output), then small heads $[64,32]$ (\code{model.py:62--101}) & Monolithic MLPs $[128,128,256,64,32] \to 2Z$ per latent, split into $(\mu, \log\sigma)$ (\code{phys\_vae.py:104--166}) --- similar capacity, different structure \\
Unmixer & Residual MLP from $(x, z_{\mathrm{aux1}}, z_{\mathrm{aux2}})$: $\mathrm{unmixed} = x + \mathrm{MLP}([x, z_1, z_2])$ (\code{model.py:96--99, 246--248}) & Learnable \code{Linear} map of $x$ $+$ residual MLP from the auxiliary \emph{means only} (no $x$ input) $+$ optional skip (\code{phys\_vae.py:139--153}) \\
In-ODE aux & $\mathrm{MLP}([z_{\mathrm{aux1}}, \theta, \dot\theta, t]) \to [64,64] \to \dot\theta'$ correction; \textbf{raw $t$} (\code{model.py:47--49}) & $\mathrm{MLP}([z, y, \sin t, \cos t]) \to [64,64] \to$ correction; $\sin/\cos$ features (\code{phys\_vae.py:238--247}) \\
Out-ODE aux & $\mathrm{MLP}([\omega, z_{\mathrm{aux1}}, z_{\mathrm{aux2}}]) \to [128,128] \to 50$-dim residual, \textbf{constant given the latents}, added to all frames including frame 0 (\code{model.py:51--53, 238--249}) & $\mathrm{MLP}([z_{\mathrm{aux2}}, \hat{x}]) \to [128,128] \to 49$-dim residual: sees the decoded series, frame 0 preserved (\code{phys\_vae.py:251, 388--403}) --- more expressive input, but loses the $\omega$ and $z_{\mathrm{aux1}}$ inputs \\
Integration & \code{odeint} Euler, \code{intg\_lev} $= 1$, $dt$ from the data file (\code{train.py:69--70}) & diffrax Euler, \code{integral\_level} $= 1$, $dt = 0.05$ derived from the dataset config --- faithful \\
Combination enumeration & Hand-written $x_P, x_{PA}, x_{PB}, x_{PAB}$ decodes & Generalized bitmask enumeration (\code{gen\_selections}, \code{full\_forward}) --- a genuine generalization \\
\bottomrule
\end{tabularx}
\end{table}

\section{Training losses and objectives}
\label{sec:losses}

Both implement the same ELBO plus three regularizers. The original's pendulum
\code{physnn} recipe (\code{pendulum\_train.sh}) versus the reimplementation's
config (\code{configs/pendulum\_vae/26MAR13.yaml}):

\begin{table}[h]
\centering
\small
\begin{tabularx}{\textwidth}{LXX}
\toprule
\textbf{Term} & \textbf{Original} & \textbf{Reimplementation} \\
\midrule
Reconstruction & $\sum$ over all 50 dims, mean over batch (\code{train.py:71}) & $\sum$ over $(T, L)$, mean over batch (\code{train\_phys\_vae.py:85}) --- match \\
KL balance & $(\code{balance\_kld} + \code{balance\_lact\_enc}) \cdot \mathrm{kld} \cdot e^{x\_\mathrm{lnvar}}$ with $x\_\mathrm{lnvar} = -9.0$, i.e.\ weight $\approx 1.01 \cdot 1.23 \cdot 10^{-4}$ (\code{train.py:196, 144}) & $\code{kl\_balance} \cdot \mathrm{kld} \cdot e^{x\_\mathrm{log\_sigma}}$ with $x\_\mathrm{log\_sigma} = 0.01$, i.e.\ weight $\approx 1.01$ (\code{train\_phys\_vae.py:237}) --- \textbf{$\sim 8{,}200\times$ stronger KL} (N1) \\
Unmixing reg.\ ($R_{\mathrm{DA},1}$) & weight $\gamma = 10^{-3}$ (\code{pendulum\_train.sh}) & $\beta = 10^{-3}$ (\code{physics\_only\_loss}, \code{train\_phys\_vae.py:94--114}) --- match \\
Data-augmentation reg.\ ($R_{\mathrm{DA},2}$) & weight $\beta = 10^{-1}$: random $\omega \sim U[0.392, 3.53]$ (the model range), physics-only decode, regress the softplus head (\code{train.py:133--140}) & $\gamma = 10^{-1}$: same construction, draws from \code{phys\_param\_ranges} $= [0.392, 3.53]$, regresses $\mu$ only (\code{train\_phys\_vae.py:115--163}) --- match, D4 fixed \\
Least-action / discrepancy ($R_\mathrm{ppc}$) & weight $\alpha = 10^{-2}$, $\tfrac14$ pairwise combination terms (\code{train.py:142--149}) & $\alpha = 10^{-2}$, per-sample Hamming-distance-1 pairs via \code{gen\_mask} (\code{train\_phys\_vae.py:199--215}) --- match, D5 fixed \\
ELBO target & $x_{PAB}$ (all auxiliary applied) & \code{x\_stat[0][-1]} (last combination) --- faithful \\
\bottomrule
\end{tabularx}
\end{table}

\noindent The pendulum recipe is numerically very different from advdif
(unmix $10^{-3}$ vs $10^{-2}$, data-augmentation $10^{-1}$ vs $10^{6}$,
discrepancy $10^{-2}$ vs $10^{-1}$), and the reimplementation transfers all
three correctly. The one large loss-side deviation is the KL balance (N1).

\section{Training protocol}
\label{sec:protocol}

\begin{table}[h]
\centering
\small
\begin{tabularx}{\textwidth}{LXX}
\toprule
\textbf{Aspect} & \textbf{Original} & \textbf{Reimplementation} \\
\midrule
Epochs & 5000 & 5000 (\code{26MAR13.yaml}) --- match \\
Optimizer & Adam, lr $10^{-3}$, $\varepsilon = 10^{-3}$, weight decay $10^{-6}$, value grad-clip 5.0 & Identical via \code{train.py} (\code{optax.clip} matches value clipping; \code{add\_decayed\_weights} matches torch L2) --- D7 inherited \\
Batch / train size & 200 / 1000 & 200 / 1000 --- match \\
Validation & No-grad evaluation; best-validation-reconstruction checkpoint saved & No-grad \code{val\_step}; best-validation checkpoint saved and restored (\code{train.py:186--216}) --- D7 fixed \\
\bottomrule
\end{tabularx}
\end{table}

\section{Evaluation}
\label{sec:eval}

\begin{table}[h]
\centering
\small
\begin{tabularx}{\textwidth}{LXX}
\toprule
\textbf{Aspect} & \textbf{Original} & \textbf{Reimplementation} \\
\midrule
Coefficient metric & Paper Table 1 reports the inference error of $\omega$ on the test split & \code{coeff\_mae}: absolute MAE of the posterior mean $\mu$ vs true $\omega$ (\code{eval.py:82}, \code{phys\_vae.py:637}) --- D6 fixed (posterior mean, not a sample), but absolute MAE is not the paper's error definition \\
Headline experiment & \textbf{Counterfactual}: scale the inferred $\omega$ by $0.2 \ldots 3.0$, decode, compare against ground-truth trajectories with scaled $\omega$ (\code{pendulum\_cntfact.ipynb}) & \textbf{No counterfactual path exists} --- \code{eval.py} reports only forecast MAE and coefficient MAE (N8) \\
Estimate used & Posterior draw in the notebook; posterior statistic for the paper's table & Posterior mean $\mu$ (\code{phys\_vae.py:637}) \\
\bottomrule
\end{tabularx}
\end{table}

\section{Ranked diagnosis of pendulum-specific deviations}
\label{sec:diagnosis}

\subsection*{N1 --- KL term effectively $8{,}200\times$ stronger (dominant)}

The original balances the KL by the observation variance $e^{-9} \approx 1.23
\cdot 10^{-4}$, which encodes the $10^{-2}$ observation noise ($\ln 10^{-4}
\approx -9.2$). The reimplementation balances by $e^{0.01} \approx 1.01$
(\code{train\_phys\_vae.py:237}, config \code{x\_log\_sigma: 0.01}). With the
KL at full strength, the posterior over $\omega$ is pulled toward the prior
--- mean $1.961$, std $2.72$ --- and the collapse regime has MAE
$\mathbb{E}|\omega - 1.961| \approx 0.59$ over the data range $[0.785, 3.14]$,
against a data range width of $2.36$. The original also adds
\code{balance\_lact\_enc} $= \alpha = 10^{-2}$ to the KL balance
(\code{train.py:196}); the reimplementation has no equivalent (a $1\%$ effect,
noted for completeness). This is the same class as the advdif D2, but the
factor is $8{,}200\times$ rather than $10^6\times$ because the original's
pendulum \code{x\_lnvar} is $-9$, not $-13.8$.

\subsection*{N2 --- No observation noise in the generated data}

The original adds $\sigma = 10^{-2}$ Gaussian noise to every frame
(\code{generate.py:88}); the reimplementation returns the raw Euler solution
(\code{pivae.py:155}). Consequences: (a) the ``correct'' \code{x\_log\_sigma}
for the original protocol would be $\approx -9.2$, so N1's fix and this item
are linked; (b) the task is easier than the original's (noiseless targets),
so forecast MAE numbers are not comparable to the original's setting; (c) the
cached dataset (\code{data/pivae/pendulum/*\_50.npy}) was generated under this
convention, so regenerating with noise invalidates the cache (the manifest
keys only on grid size, \code{x\_grid}, and \code{time}, not on the noise
convention).

\subsection*{N3 --- Euler-generated data versus exact trajectories}

DOP853 versus Euler at $dt = 0.05$ with $\omega \le \pi$: the generated
trajectories deviate from the true solutions of the same ODE (noticeably for
the high-$\omega$, high-$A$ samples). As noted in Section~\ref{sec:data}, the
model's Euler physics makes the decomposition self-consistent, so this does not
break identifiability --- but any evaluation against externally computed
``true'' trajectories (e.g.\ the original's counterfactual ground truth) would
need to use the same Euler discretization to be meaningful.

\subsection*{N4 --- Linear $\omega$ head with hard clipping}

The original's softplus head guarantees $\omega > 0$ smoothly and then clamps
to a wide feasible range $[0, 20]$; the reimplementation's linear $\mu$ can go
negative and the sampled latent is hard-clipped to $[0.392, 3.53]$
(\code{phys\_vae.py:748}), with zero gradient outside. Because the model range
strictly contains the data range $[0.785, 3.14]$, this is far less damaging
than the advdif variant of the same issue (where the clip equalled the data
range) --- but early in training, before the encoder scales, a substantial
fraction of samples can sit at the clip boundary with dead gradients.

\subsection*{N5 --- $\sin t, \cos t$ features instead of raw $t$ in the in-ODE auxiliary}

The forcing term the pre-function must absorb is
$A\omega^2\cos(2\pi\phi t)$ with $\phi \in [\pi, 2\pi]$: $1.2$--$2.4$
oscillations over the $2.45$\,s span. Neither raw $t$ (original) nor
$(\sin t, \cos t)$ (reimplementation) expresses these frequencies directly ---
the network must synthesize them from $z$ and the state $(\theta, \dot\theta)$
--- but raw $t$ is an unbounded, monotone feature over the span, while
$(\sin t, \cos t)$ is a bounded, frequency-1 basis. Both models share the
state inputs, so this is a capacity nuance rather than a bug; it is listed
because the pendulum is the case where the auxiliary's job (arbitrary-frequency
forcing) is hardest.

\subsection*{N6 --- Unmixer residual does not see the data}

The original's unmixer residual MLP takes $(x, z_{\mathrm{aux1}},
z_{\mathrm{aux2}})$ as input; the reimplementation's residual MLP takes only
the auxiliary means, with the data entering through a separate learnable
\code{Linear} map and optional skip (\code{phys\_vae.py:139--153}). The
reimplementation therefore cannot make the unmixing correction
\emph{locally} dependent on the observed series, only globally through a
linear map. Also, the config key \code{unmixer.latent\_dim: 50} is silently
ignored (the input width is the sum of auxiliary latent dims).

\subsection*{N7 --- Out-ODE auxiliary input/output differences}

The original's out-ODE residual is a function of $(\omega, z_{\mathrm{aux1}},
z_{\mathrm{aux2}})$ only --- constant across time given the latents --- and is
added to every frame including frame 0. The reimplementation's post-function
sees the decoded series $\hat{x}$ (more expressive) but not $\omega$ or
$z_{\mathrm{aux1}}$, and never modifies frame 0. For the pendulum, the
out-ODE auxiliary's residual job (whatever the in-ODE auxiliary cannot fix)
is modest, so this is a design divergence to be aware of rather than a
ranked error.

\subsection*{N8 --- No counterfactual evaluation path}

The original's headline pendulum result is the counterfactual experiment:
rescale the inferred $\omega$ and decode. The reimplementation's evaluation
(\code{eval.py}) computes only forecast and coefficient MAE; there is no API
to decode from arbitrary $(\omega, z)$ with a fixed initial condition. Since
the counterfactual is the paper's evidence that $\omega$ is
\emph{interpretable} (not merely estimable), the reimplementation currently
cannot reproduce the experiment that motivates the architecture.

\subsection*{N9 --- Encoder structure}

Monolithic per-latent MLPs versus feature extractor + small heads. Similar
parameter counts, no identified failure mode; listed for completeness.

\section{PAID on the pendulum: what must change}
\label{sec:paid}

PAID (\code{src/models/paid.py}) shares the physics model, the dataset, the
combination machinery, and the trainer with PhysVAE, and the earlier P1/P3/P5
fixes (post-function latent slice, coefficient head range, reconstruction
over the noised region only) are present in the current code. The regularizer
weights in the committed pendulum configs (\code{disc\_reg} $10^{-2}$,
\code{phys\_reg} $10^{-3}$, \code{aug\_reg} $10^{-1}$) match the original
\code{physnn} recipe exactly --- unlike the advdif case, no weight audit is
forced by the original's pendulum recipe. The items below are the
pendulum-specific gaps, ranked by impact.

\subsection*{A1 --- Second-order ODE state initialization at the condition frame (structural)}

PAID's denoiser solves the physics ODE starting at the last condition frame:
\code{t0 = t\_span[condition-1]}, with the initial state taken from
\code{x[:, condition-1]} and \textbf{the derivative components of the state
padded with zeros} (\code{paid.py:67--68, 320--327}). For first-order systems
(advection--diffusion, Burgers) the padding is vacuous. The pendulum is
2-D state $(\theta, \dot\theta)$: the solve starts with
$\dot\theta(t_0) = 0$ when the true angular velocity at that frame is not
zero. Measured on generated data (central difference at frame 4,
\code{condition} $= 5$, $dt = 0.05$): mean $|\dot\theta| \approx 0.59$, max
$3.85$, and $80\%$ of samples exceed $0.1$.

Consequences: the physics-only combination (pre/post functions off) --- which
anchors the discrepancy loss and embodies the ``known physics'' claim ---
starts from a wrong velocity and is systematically wrong no matter how well
$\omega$ is estimated; the discrepancy loss then measures this initialization
error rather than the auxiliary functions' deviation from physics; and the
pre-function is forced to spend capacity steering $\dot\theta$ from zero
instead of modelling damping/forcing. Note the asymmetry inside PAID itself:
\code{generate\_physics} (the \code{phys\_reg} target) solves from frame 0
with $\dot\theta(0) = 0$, which is correct \emph{for this dataset}
(\code{train\_paid.py:149--150}), so the unmixing target and the denoiser's
internal solve disagree on their initial state. PhysVAE and the original both
decode from frame 0 and are unaffected.

This is the fundamental change PAID needs for the pendulum: the derivative
initial condition at the solve anchor must come from somewhere other than a
zero pad --- e.g.\ estimated from the condition frames (they are clean, so a
finite difference or a learned readout is available), made a model output, or
the solve anchored at frame 0 with the condition used only for guidance.
The choice is a design decision; the observation is that the current padding
is wrong for $80\%$ of pendulum samples.

\subsection*{A2 --- Committed pendulum PAID configs do not construct}

Verified by execution: loading \code{configs/pendulum\_paid/base.yaml} and
calling the constructor fails with \code{TypeError: PAID.\_\_init\_\_() got an
unexpected keyword argument 'unet'}. The same class of staleness as the
earlier P2, still present for the pendulum configs: a dead \code{unet:} block;
\code{features}/\code{in\_features}-style auxiliary/physics functions where
the current \code{make\_mlp} interface expects \code{dims};
\code{integral\_level: 1.0} as a float; \code{condition\_base.yaml} carries
the literal \code{[CONDITION]} placeholder, legacy \code{pde:} keys
(\code{system}, \code{timestep}, \code{spatial\_length},
\code{num\_grid\_points}, \code{time\_dim}) that the constructor rejects, and
a \code{trainer\_config} without \code{weight\_decay}/\code{grad\_clip}, so
the inherited D7 protocol is silently inactive. A hand-written corrected
config (dims-style functions, no \code{unet} block, \code{condition: 5})
constructs and runs \code{paid\_loss} and the DDIM sampling call end-to-end
on pendulum-shaped data, so the model itself is pendulum-ready; only the
committed configs are not.

\subsection*{A3 --- Forcing representation in the pre-function}

The same $(\sin t, \cos t)$ basis as N5 applies to PAID's pre-functions
(\code{paid.py:296--310}), and the pendulum's forcing
$A\omega^2\cos(2\pi\phi t)$ is the hardest signal an in-ODE auxiliary can be
asked to carry. PAID additionally interleaves the ODE solve with the diffusion
process (the solve runs at every diffusion step), so any error the pre-function
makes is re-injected at every step. Whether the latent $z$ (dim 1 in the
committed config) plus the state inputs suffice to encode $(A, \phi)$ is an
open capacity question for the pendulum.

\subsection*{A4 --- Out-ODE auxiliary design}

The committed pendulum configs give the post-function \code{latent\_dim: 1},
where the original's \code{physnn} recipe uses $z_{\mathrm{aux2}} \in
\mathbb{R}^2$ (and the reimplementation's PhysVAE pendulum config follows the
original). PAID's post-function input is $[z, \hat{x}]$ --- like PhysVAE's
reimplementation it does not receive $\omega$ or the pre-function's latent.
For the pendulum the out-ODE auxiliary is the last line of defence for
whatever the in-ODE auxiliary cannot represent; whether one latent dimension
suffices is untested.

\subsection*{A5 --- Training budget and batch size}

The committed pendulum PAID configs use batch 64 and 100 epochs: $100 \times
\lceil 1000/64 \rceil = 1{,}600$ gradient steps, versus the original's
$5{,}000 \times 5 = 25{,}000$ steps at batch 200. PAID's per-step cost is
higher (an $m$-combination ODE solve per diffusion step), but a $16\times$
smaller step count with a $3\times$ smaller batch is a substantial reduction
in exposure to the data, and unlike PhysVAE there is no KL to keep the
latents honest --- the coefficient estimate is supervised only by
\code{aug\_reg} ($10^{-1}$, matching the original's pendulum recipe).

\subsection*{A6 --- Coefficient read-off protocol}

At evaluation, PAID noises the generated sample back by a single diffusion
step and reads the embedder's physics parameters at $\tau = 1$
(\code{paid.py:645--656}). With the committed linear schedule
($\beta_1 = 10^{-4}$), $\alpha_1 \approx 0.9999$: the embedder sees an
essentially clean sample, so the protocol is close to ``read the coefficient
from the generated series''. For the pendulum this is mechanically sound
(verified in the smoke test), but the estimate is a byproduct of the
generated trajectory rather than an inference from the observed condition ---
the same protocol question flagged as P6, unresolved.

\subsection*{A7 --- Shared dataset deviations}

PAID trains on the same generated pendulum dataset as PhysVAE, so N2 (no
observation noise) and N3 (Euler trajectories) apply unchanged. For PAID the
noiselessness has a different flavour: diffusion training adds its own noise
process, so the observation-noise question is entangled with the choice of
$\beta$ schedule rather than with a KL balance. Nothing in the committed
schedule ($\beta \in [10^{-4}, 0.02]$, 1000 steps) references the data's
noise scale; the original's $\sigma = 10^{-2}$ is the natural calibration
point if that correspondence is intended.

\section{What was and was not verified by execution}

\begin{itemize}[nosep]
  \item Verified: \code{configs/pendulum\_vae/26MAR13.yaml} constructs;
        \code{phys\_vae\_loss} and the \code{PhysVAE.\_\_call\_\_} reconstruction
        path run on generated $(8, 50, 1)$ pendulum data;
        \code{configs/pendulum\_paid/base.yaml} fails to construct
        (\code{TypeError: 'unet'}); a corrected PAID config constructs and runs
        \code{paid\_loss} and the DDIM sampling call on the same data;
        $|\dot\theta|$ at the condition frame has mean $0.59$, max $3.85$
        (central differences on generated data).
  \item Not verified: any full training run, coefficient-estimation accuracy,
        the empirical size of N1's KL-collapse pressure on the pendulum
        (the reconstruction term may win against the KL here --- the pendulum
        series is more sensitive to $\omega$ than advdif is to $a$), and all
        paper-number comparisons. The ranking of N1--N3 is from code analysis;
        N1 is confirmable in one short run by checking whether the predicted
        $\omega$ clusters near the prior mean $1.961$.
\end{itemize}

\end{document}
